\section*{第 \ref{ch:g8:reproductive-systems} 章——生殖系统}

\begin{solution}{exo:g8:reproductive-systems:1}
阴囊中的\omterm{def:g8:reproductive-systems:male}{睾丸}从\omterm{def:g8:puberty:puberty}{青春期}起
持续产生\omterm{def:g8:reproductive-systems:male}{精子}。输精管道和腺体：
负责储存、运输，并提供构成精液的液体。\omterm{def:g8:reproductive-systems:male}{阴茎}：
是精液的排出通道，也是\omterm{def:g7:removing-wastes:kidneys}{尿液}的排出通道，两者由\omterm{def:g7:heart-and-circulation:rooms}{瓣膜}分工把关。
\end{solution}

\begin{solution}{exo:g8:reproductive-systems:2}
\omterm{def:g8:reproductive-systems:female}{卵巢}：自出生起储备\omterm{def:g4:animal-reproduction:sexual}{卵细胞}，
从\omterm{def:g8:puberty:puberty}{青春期}起每个周期排出一个。\omterm{def:g8:reproductive-systems:female}{输卵管}：承接排出的\omterm{def:g4:animal-reproduction:sexual}{卵细胞}，
是\omterm{def:g5:flower-to-fruit:steps}{受精}的常见场所。\omterm{prop:g5:human-reproduction:plan}{子宫}：\omterm{def:g3:bones-and-muscles:muscle}{肌肉}质的
容纳\omterm{def:g5:human-reproduction:pregnancy}{胚胎}的器官，内膜定期更新。阴道：接纳
精液，并作为产道。
\end{solution}

\begin{solution}{exo:g8:reproductive-systems:3}
\omterm{def:g8:reproductive-systems:male}{精子}：人体最小的\omterm{def:g5:first-look-at-cells:cell}{细胞}，头部含\omterm{def:g6:cell-unit-of-life:parts}{细胞核}，
尾部用于游动，储备物质少，
由数亿个持续不断地产生。\omterm{def:g4:animal-reproduction:sexual}{卵细胞}：人体最大的\omterm{def:g5:first-look-at-cells:cell}{细胞}，
圆形、不动，储备丰富，
每月释放一个。
\end{solution}

\begin{solution}{exo:g8:reproductive-systems:4}
\omterm{def:g8:reproductive-systems:female}{卵巢}每次排出一个\omterm{def:g4:animal-reproduction:sexual}{卵细胞}，由\omterm{def:g8:reproductive-systems:female}{输卵管}接纳，
约在第 14 天；\omterm{def:g5:flower-to-fruit:steps}{受精}能力维持约一天。
\end{solution}

\begin{solution}{exo:g8:reproductive-systems:5}
\omterm{prop:g5:human-reproduction:plan}{子宫}准备好的内膜逐渐分解、脱落，
持续数天。这说明本周期
未发生\omterm{def:g5:flower-to-fruit:steps}{受精}。
\end{solution}

\begin{solution}{exo:g8:reproductive-systems:6}
\omterm{def:g5:flower-to-fruit:steps}{受精}。\omterm{prop:g8:sexual-reproduction:core}{受精卵}着床于已准备好的内膜，
周期暂停，\omterm{def:g5:human-reproduction:pregnancy}{怀孕}开始。
\end{solution}

\begin{solution}{exo:g8:reproductive-systems:7}
第一步，\omterm{def:g8:sexual-reproduction:gametes}{配子}：\omterm{def:g8:reproductive-systems:male}{睾丸}和\omterm{def:g8:reproductive-systems:female}{卵巢}。第二步，\omterm{def:g5:flower-to-fruit:steps}{受精}：
在\omterm{def:g8:reproductive-systems:female}{输卵管}中发生，\omterm{def:g8:reproductive-systems:male}{精子}由阴道经\omterm{prop:g5:human-reproduction:plan}{子宫}游至此处。
第三步，\omterm{prop:g8:sexual-reproduction:core}{受精卵}的分裂发育：
着床于\omterm{prop:g5:human-reproduction:plan}{子宫}已准备好的内膜中。
\end{solution}

\begin{solution}{exo:g8:reproductive-systems:8}
男性一方遵循“数量多、成本低”的策略：持续
产生数以百万计的\omterm{def:g8:reproductive-systems:male}{精子}，每个都可以被舍弃。
女性一方遵循“数量少、储备丰富”的策略：每月
只有一个储备充足的\omterm{def:g4:animal-reproduction:sexual}{卵细胞}，背后则是\omterm{prop:g5:human-reproduction:plan}{子宫}为容纳\omterm{def:g5:human-reproduction:pregnancy}{胚胎}所做的全部投入。
这两种策略，都写在\omterm{def:g5:first-look-at-cells:cell}{细胞}之中。
\end{solution}

\begin{solution}{exo:g8:reproductive-systems:9}
这套机制在最初几年仍在调整，\omterm{def:g8:puberty:hormones}{激素}
链条尚未稳定，因此月经不规律是正常的；
周期长短确实因人而异，也因月而异。
这两者都表明系统运行正常。
\end{solution}

\begin{solution}{exo:g8:reproductive-systems:10}
\omterm{prop:g5:human-reproduction:plan}{子宫}是一块\omterm{def:g3:bones-and-muscles:muscle}{肌肉}，脱落本身是一种\omterm{def:g3:bones-and-muscles:muscle}{肌肉}活动——
其收缩即表现为痉挛。温暖和活动
可以缓解（若效果不够，
校医会提供建议）。
\end{solution}

\begin{solution}{exo:g8:reproductive-systems:11}
\omterm{def:g8:reproductive-systems:male}{精子}的任务是“远行”：成本低、能游动，
因几乎无法全部到达而需数以百万计，所以生产
持续不断。\omterm{def:g4:animal-reproduction:sexual}{卵细胞}的任务是成为\omterm{prop:g8:sexual-reproduction:core}{受精卵}最初的
营养来源：成本高，自出生即已储备，
每月排出时须经历一个月的准备，因此它的“工厂”
更像一座宝库，每次只释放一个。
\end{solution}

\begin{solution}{exo:g8:reproductive-systems:12}
第一阶段是准备待解的问题：内膜重建，
容纳\omterm{def:g6:exploring-our-environment:conditions}{条件}就绪。第二阶段是提出问题：
排出一个\omterm{def:g4:animal-reproduction:sexual}{卵细胞}——能否\omterm{def:g5:flower-to-fruit:steps}{受精}？第三阶段是“否”的答案：
无用的内膜脱落，周期于下月再次发问。第四阶段是“是”的答案：
\omterm{prop:g8:sexual-reproduction:core}{受精卵}着床，发问停止，\omterm{prop:g5:human-reproduction:plan}{子宫}
作为容纳之所开始发挥作用，
持续九个月。
\end{solution}

\begin{solution}{pb:g8:reproductive-systems:1}
\textbf{1.} \omterm{def:g8:reproductive-systems:male}{精子}产生于\omterm{def:g8:reproductive-systems:male}{睾丸}，即阴囊中的“工厂”，
进入输精管道储存和运输，
与腺体分泌的液体混合成精液，
经\omterm{def:g8:reproductive-systems:male}{阴茎}排出；同一器官在其他时候排出\omterm{def:g7:removing-wastes:kidneys}{尿液}，
由\omterm{def:g7:heart-and-circulation:rooms}{瓣膜}安排，使两种功能
严格分开。

\textbf{2.} \omterm{def:g4:animal-reproduction:sexual}{卵细胞}储备于\omterm{def:g8:reproductive-systems:female}{卵巢}中，自出生起即已具备；
\omterm{def:g5:flower-to-fruit:steps}{受精}通常发生在\omterm{def:g8:reproductive-systems:female}{输卵管}中；\omterm{def:g5:human-reproduction:pregnancy}{怀孕}时
\omterm{prop:g8:sexual-reproduction:core}{受精卵}着床于\omterm{prop:g5:human-reproduction:plan}{子宫}已准备好的内膜中。

\textbf{3.} 两种\omterm{def:g8:sexual-reproduction:gametes}{配子}的构造遵循相反的策略：
数以百万计的廉价游动者，对应每月一个储备充足的\omterm{def:g4:animal-reproduction:sexual}{卵细胞}——
数量与投入的对立，
写在\omterm{def:g5:first-look-at-cells:cell}{细胞}本身之中。

\textbf{4.} 系统早已建成，只是处于闲置，
等待脑底部调控中心的\omterm{def:g8:puberty:hormones}{激素}，
于\omterm{def:g8:puberty:puberty}{青春期}将其启动。

\textbf{5.} 大约第 1--5 天：\omterm{prop:g8:reproductive-systems:cycle}{月经期}。直到周期中段：
\omterm{prop:g5:human-reproduction:plan}{子宫}内膜重建。第 14 天前后：\omterm{prop:g8:reproductive-systems:cycle}{排卵}，
易受孕日位于其前后。然后是准备就绪的等待——直到
第 28 天，车轮开始下一轮转动。

\textbf{6.} 周期长短因人而异，
同一个人的周期逐月也不相同——而头几年，
机制仍在调整，周期往往不规则。日历所呈现的是
规律，而非任何人确切的日期。

\textbf{7.} 取消：内膜剥落与下一次
\omterm{prop:g8:reproductive-systems:cycle}{月经期}——车轮停转。开始：\omterm{prop:g8:sexual-reproduction:core}{受精卵}着床于内膜，
以及\omterm{def:g5:human-reproduction:pregnancy}{怀孕}的九个月。

\textbf{8.} 痉挛是\omterm{prop:g5:human-reproduction:plan}{子宫}\omterm{def:g3:bones-and-muscles:muscle}{肌肉}在脱落时的收缩——
属正常现象，并非损伤；保暖与活动有所帮助，
若疼痛更为剧烈，应去
护士或医生处就诊。

\textbf{9.} 例如：“这些是像\omterm{def:g4:heart-and-blood:heart}{心脏}和肺一样的器官——
在座的每个人都拥有其中一套，
而这一章讲的正是它们。”

\textbf{10.} 工厂：\omterm{def:g8:reproductive-systems:male}{睾丸}——\omterm{def:g8:reproductive-systems:female}{卵巢}。\omterm{def:g8:sexual-reproduction:gametes}{配子}：
持续游动的数百万个——每月一个
备足养分的\omterm{def:g4:animal-reproduction:sexual}{卵细胞}。相遇：\omterm{def:g8:reproductive-systems:female}{输卵管}。住所：无——
是每月准备好的\omterm{prop:g5:human-reproduction:plan}{子宫}。时间表：从\omterm{def:g8:puberty:puberty}{青春期}起持续——
周期性的，从\omterm{def:g8:puberty:puberty}{青春期}到约五十岁。

\textbf{11.} 两个模型上的第 1 步：工厂制造
\omterm{def:g8:sexual-reproduction:gametes}{配子}。女性模型上的第 2 步：二者
在\omterm{def:g8:reproductive-systems:female}{输卵管}中结合。模型与日历贯穿的第 3 步：\omterm{prop:g8:sexual-reproduction:core}{受精卵}
在周期日历所示已准备好的内膜中着床，车轮
停转，发育开始。

\textbf{12.} “因为身体更依赖知识而非
传说——而今天所讲的每一项内容，都有一个训练有素的
保密去处，以解答后续的问题。”
\end{solution}
